\section{HSM Selection Matrix}

\label{sec:hsm}
%% =====================================================================

The five cold-tier and three warm-tier shares require eight HSM
endpoints across at least four distinct vendor lines. This section
documents each chosen device, its certification, its public price
basis, and the reason it was selected.

\subsection{Selection Principles}

\begin{enumerate}[nosep]
\item \textbf{Vendor diversity.} No two shares of the same tier may
share a single firmware codebase. AWS CloudHSM and GCP Cloud HSM both
use Marvell LiquidSecurity firmware and therefore each tier admits at
most one Marvell-backed share.
\item \textbf{Certification floor.} FIPS 140-2 Level 3 minimum on
every signing device. FIPS 140-3 acceptable as a substitute where
the validation is current. Level 4 physical (Utimaco SR-Se200) used
in cold tier for the on-prem DR share.
\item \textbf{Operational reachability.} On-prem HSMs require redundant
power, network, and physical access control matching SOC 2 CC6.4 and
CCSS data-centre controls.
\item \textbf{Supply chain.} The fund retains acceptance authority
over firmware updates: no signer auto-updates from vendor cloud
without a documented internal review.
\end{enumerate}

\subsection{Device-by-Device Specification}

\paragraph{Thales Luna Network HSM 7 -- Share $\mathbf{c}_1$.}
Network-attached PCIe-class HSM. FIPS 140-2 Level 3 certified
(certificate \#3205) \cite{thales-luna7}. Performance: 20{,}000
ECDSA secp256k1 signatures/sec. Public list price approximately
\$15{,}000 per appliance for the base ECDSA-capable SKU; HA pair plus
PED authentication tokens add roughly \$8{,}000. Selected as the
production-grade industry baseline; widely audited and accepted by
banking compliance teams.

\paragraph{Utimaco u.trust SecurityServer SR-Se200 Gen2 -- Share $\mathbf{c}_2$.}
PCIe Gen2 HSM with FIPS 140-2 Level 3 certification and Level 4
physical security (active mesh, environmental sensors, instant zeroisation
on tamper) \cite{utimaco-secsr}. Throughput: 5{,}500 ECDSA/sec.
Price: approximately \$25{,}000 per unit; supports custom firmware
modules signed by the operator. Selected for the DR site specifically
because the Level 4 physical envelope provides resistance to
sustained physical attack (drilling, X-ray, decapping) at a level
that the Marvell and Thales L3 devices do not match.

\paragraph{AWS CloudHSM (Marvell LiquidSecurity) -- Share $\mathbf{c}_3$.}
Single-tenant cloud HSM cluster, FIPS 140-2 Level 3 certificate
\#4218 \cite{aws-cloudhsm}. Marvell LiquidSecurity firmware. Pricing
is hourly: approximately \$1.45/hour per HSM in us-east-1, two HSMs
minimum for HA, total approximately \$2{,}100/month per cluster.
Cold-tier requires multi-region (us-east-1 + eu-west-1) so two
clusters: approximately \$4{,}200/month.

\paragraph{GCP Cloud HSM (Marvell LiquidSecurity) -- Share $\mathbf{c}_4$.}
Multi-tenant cloud HSM service backed by Marvell LiquidSecurity
modules; FIPS 140-2 Level 3 attested per Google's published security
overview \cite{gcp-cloudhsm}. Pricing per signing operation:
approximately \$0.03 per 10{,}000 ops plus \$1/hour key holding.
Cold-tier signing volume is dominated by ceremony scheduling, not
ops, so the steady-state cost is the key-holding fee: approximately
\$1{,}460/month for two regions (us-central1 + europe-west1).

\paragraph{YubiHSM 2 -- Share $\mathbf{c}_5$.}
Nano-form HSM, USB-A, FIPS 140-2 Level 3 certificate \#3915
\cite{yubihsm2}. Price: approximately \$650 per unit; spare units
maintained at the vault. Selected for the airgap share specifically
because its small size and low power envelope make it amenable to
deployment in a fully offline vault, transport in tamper-evident
bags, and physical destruction at end-of-life.

\paragraph{Azure Dedicated HSM (Thales Luna 7) -- Share $\mathbf{w}_3$.}
Dedicated single-tenant Luna 7 in an Azure datacentre, managed by
the customer \cite{azure-dedicated-hsm}. FIPS 140-2 Level 3.
Pricing: approximately \$3.80/hour per device, \$2{,}750/month per
device. Selected for warm tier as the third independent cloud
jurisdiction and to provide a Thales-line presence in the warm
quorum (defence-in-depth against a Marvell-firmware-line vulnerability
affecting both AWS and GCP shares simultaneously).

\subsection{Selection Matrix}

\begin{table}[h]
\centering
\small
\begin{tabular}{llllll}
\toprule
\textbf{Device} & \textbf{Cert} & \textbf{Vendor} &
\textbf{Throughput} & \textbf{Capex} & \textbf{Opex} \\
\midrule
Thales Luna 7        & 140-2 L3      & Thales  & 20k/s
                     & \$15k & --                              \\
Utimaco SR-Se200     & 140-2 L3+L4ph & Utimaco & 5.5k/s
                     & \$25k & --                              \\
AWS CloudHSM         & 140-2 L3      & Marvell & 1k/s
                     & --     & \$4.2k/mo                       \\
GCP Cloud HSM        & 140-2 L3      & Marvell & shared
                     & --     & \$1.5k/mo                       \\
YubiHSM 2            & 140-2 L3      & Yubico  & 30/s
                     & \$650  & --                              \\
Azure Dedicated HSM  & 140-2 L3      & Thales  & 20k/s
                     & --     & \$2.8k/mo                       \\
\bottomrule
\end{tabular}
\caption{HSM selection matrix. Capex per appliance; opex per
month per cluster. Cloud-HSM throughput differs because per-tenant
slicing is service-defined, not device-defined.}
\label{tab:hsm-matrix}
\end{table}

\subsection{What Was Excluded and Why}

\begin{itemize}[nosep]
\item \textbf{Single-vendor stack.} Excluded; one firmware vulnerability
across the entire estate is not survivable at this AUM.
\item \textbf{Software-only TEE (Intel SGX, AMD SEV-SNP).} Excluded
from the cold tier; SGX and SEV have a recurring history of
microarchitectural side-channel vulnerabilities. Acceptable as a
secondary attestation surface but not as a root signing device.
\item \textbf{Multi-tenant cloud KMS at FIPS 140-2 Level 2.} Excluded
from cold and warm tiers; admitted only at the hot tier (Tier 3) as
shares behind a 24-hour time-lock with a per-day spend cap.
\item \textbf{Smartcards as cold-share devices.} Excluded; smartcard
throughput, durability, and battery-backed RAM lifetimes do not match
multi-decade custody horizons.
\end{itemize}

\subsection{Firmware-Update Policy}

The custody operations team retains a firmware-update review queue.
A vendor-published firmware release is evaluated against three
criteria before acceptance:

\begin{enumerate}[nosep]
\item \textbf{Independent CVE review} of the release notes by the
fund's security team.
\item \textbf{Staged rollout}: deployed first to one cold share and
observed for fourteen calendar days before propagation.
\item \textbf{Cryptographic acceptance test}: a known-answer test
suite is run against the upgraded device using a published test
vector set; mismatches halt the rollout.
\end{enumerate}

This procedure intentionally accepts that the fund will run firmware
slightly behind the vendor's latest, in exchange for the time needed
to verify that an upgrade does not silently change signing behaviour.

%% =====================================================================
